\documentclass[11pt]{article}
\usepackage[T1]{fontenc}
\usepackage{lmodern,amsmath,amssymb,amsthm,mathtools,microtype}
\usepackage[margin=1.04in]{geometry}
\usepackage[colorlinks=true,linkcolor=blue!45!black,urlcolor=blue!45!black,citecolor=blue!45!black]{hyperref}
\usepackage{xcolor}
\newtheorem{theorem}{Theorem}[section]
\newtheorem{lemma}[theorem]{Lemma}
\newtheorem{corollary}[theorem]{Corollary}
\theoremstyle{definition}\newtheorem{example}[theorem]{Example}
\newcommand{\OO}{\mathcal O}
\newcommand{\Cp}{\mathbb C_p}
\newcommand{\Zp}{\mathbb Z_p}
\newcommand{\U}{\mathcal U}
\newcommand{\ord}{\operatorname{ord}}
\newcommand{\ceil}[1]{\left\lceil #1\right\rceil}
\title{The exact cost of simultaneous nonvanishing\\at powers of a torsion point}
\author{A research note accompanying the ProveIt analysis}
\date{27 November 2026}
\begin{document}
\maketitle
\begin{abstract}
On a primitive cyclotomic shell, the least valuation budget that can force a
zero at one of a prescribed set of unit powers equals a finite group covering
density. The formula is exact in every dimension, for arbitrarily many bounded
integral power series with separate witnesses. It follows by forming a product
over the unit group and counting its order of vanishing on a Cartesian grid.
For the first two powers at an odd prime, we determine the optimal endpoint
over consecutive shells in every dimension at every half-integral budget.
Whether one saves a shell is decided by the parity of the multiplicative
order of $2$ modulo $p$. The exact cost on any collection of shells is also
determined in one variable. All arguments are local and unconditional.
\end{abstract}

\section{The selection problem}
A twisting step in the source development asks for one torsion tuple for which
several power specializations satisfy valuation inequalities. Its nonvanishing
version is isolated in the ProveIt article, section on torsion specialization,
Proposition \emph{Selection for the first $m$ powers}, part (a) \cite{proveit}.
The question \emph{Several powers} asks when that proposition's endpoint is
optimal. We address nonvanishing for powers relatively prime to $p$.
The first two powers at odd primes are already sensitive to the arithmetic
of $p$. Section~\ref{sec:two} gives exact endpoints in the article's
normalization.

Write $v(p)=1$ on $\Cp$, put $v(0)=+\infty$, and write $\OO$ for the valuation
ring. At level $h\ge1$ let
\[
 e_h=(p-1)p^{h-1},\qquad r_h=1/e_h,\qquad
 Z_h=\{\zeta-1:\ord(\zeta)=p^h\},\qquad
 \U_h=(\mathbb Z/p^h\mathbb Z)^\times .
\]
The shell consists of $e_h$ distinct points of valuation $r_h$, by the
Eisenstein polynomial $g_h(T)=\Phi_{p^h}(1+T)$. For nonempty finite $S$ put
$A=\bigcup_{h\in S}Z_h$, $H=\max S$, and $N=|S|$. Fix $s\ge1$, $r\ge\max_{h\in S}r_h$
and a finite nonempty set $J$ of positive integers prime to $p$.
An admissible family is a finite collection of $F_i\in\OO[[T_1,\ldots,T_s]]$ and witnesses
$x_i\in\Cp^s$ with every $v(x_{i,b})\ge r$ and finite total budget
$B=\sum_i v(F_i(x_i))$. A family \emph{blocks} $A^s$ if for each $z\in A^s$
there are $i$ and $j\in J$ with $F_i([j]z)=0$, where
$[j]z=((1+z_b)^j-1)_{b=1}^s$. Write $C_s(S,J)$ for the infimum of the
budgets of blocking families, with $r$ fixed. The formulas below hold for
every permitted $r$. All upper bounds are attained by one polynomial in
one coordinate with common witness zero.

Let $J_h$ be the image of $J$ in $\U_h$. Define
\[
 \tau_h=\min\bigl\{|D|:\ D\subseteq\U_h,\ 
        (uJ_h)\cap D\ne\varnothing\ \text{for every }u\in\U_h\bigr\},
 \qquad c_h=\frac{\tau_h}{e_h}.
\]
Thus $D J_h^{-1}=\U_h$. This is an ordinary finite covering number. Exact
factorizations of finite groups and the related $k$-radius prime problem
have an extensive independent literature \cite{blackburn}.
Our use of the covering density identifies an extremum for analytic
functions of arbitrarily many variables.

\begin{theorem}\label{th:cover}
For a single shell, $C_s(\{h\},J)=c_h$ in every dimension. For any $S$,
\[
 C_1(S,J)=\sum_{h\in S} c_h,\qquad
 \max_{t\in S}\bigl(|S\cap[1,t]|\,c_t\bigr)\ \le C_s(S,J)\le
 \sum_{h\in S}c_h .
\]
In particular the exact value is $Nc$ in every dimension if all the $c_h$
on $S$ equal $c$. For $B<C_s$ every admissible family of budget at most $B$
admits a tuple at which all powers are nonzero. In every exact formula above
a blocker exists at equality.
\end{theorem}

\section{A grid estimate with multiplicities}
For an integral series $P=\sum a_\nu T^\nu$, write
$G_r(P)=\min_\nu(v(a_\nu)+r|\nu|)$, with $G_r(0)=+\infty$.
The Gauss valuation is additive on products: for two nonzero series one
may take, in a fixed monomial order, the least index attaining each minimum;
their sum has a unique contributing pair among terms of minimal weight at
that index. Also $v(P(x))\ge G_r(P)$ on the witness polydisc.
Substitution by $[u]$ for $p\nmid u$ preserves $G_r$. Indeed both the forward
series and its compositional inverse $(1+T)^{1/u}-1$ have integral
coefficients and zero constant term. Substitution by either therefore cannot
decrease $G_r$. The inverse is integral since $1/u\in\Zp$ and binomial
coefficients are continuous and integral on $\Zp$.

Ball and Serra proved the multiplicity grid-ideal identity for polynomials
over a field \cite[Theorem~3.1]{ballserra}. Here integral quotients for
bounded series are needed to retain the Gauss bound.
\begin{lemma}\label{lem:grid}
Let $g=\prod_{a\in A}(T-a)\in\Zp[T]$ be monic with distinct roots in the
open unit disc and $g\equiv T^d\pmod p$, $d=\deg g$. Suppose an integral
series $P$ vanishes to order at least $M\ge1$ at each point of $A^s$
(all local terms of total degree below $M$ vanish). Then
\[
 P\in(g(T_1),\ldots,g(T_s))^M\OO[[T_1,\ldots,T_s]],
 \qquad G_r(P)\ge M G_r(g).
\]
\end{lemma}
\begin{proof}
Divide successively in each variable by $g(T_b)^M$ to write
$P=\sum_b g(T_b)^M Q_b+R$ with integral quotients and $R$ of degree
below $Md$ in each variable. For completeness, for a monic divisor
$T^l+b(T)$ with all coefficients of $b$ divisible by $p$, let $S_l$ discard
the first $l$ coefficient slices and shift the others down. The quotient is
$(1+S_l b)^{-1}S_lP$, where $b$ acts by multiplication. The geometric
series for this inverse converges in the supremum norm of bounded coefficient
arrays. This also preserves previous degree bounds in other variables.

Expand the polynomial remainder uniquely by repeated monic division as
\[
 R=\sum_{0\le\beta_b<M} S_\beta(T)\prod_b g(T_b)^{\beta_b},
 \qquad \deg_{T_b}S_\beta<d .
\]
It still has multiplicity at least $M$ throughout $A^s$.
Induct on $k<M$ to show $S_\beta=0$ for $|\beta|=k$. At $a\in A^s$,
after all smaller terms have disappeared, the homogeneous term of degree $k$
in the increments $X_b=T_b-a_b$ is
$\sum_{|\beta|=k}S_\beta(a)\prod_b(g'(a_b)X_b)^{\beta_b}$.
Distinct monomials and simple roots force all $S_\beta(a)=0$.
Interpolation in one coordinate at a time gives $S_\beta=0$ identically.
The surviving terms belong to the asserted power of the ideal.
All multipliers are integral, so its Gauss bound follows.
\end{proof}

\begin{proof}[Proof of Theorem~\ref{th:cover}]
Put $F=\prod_i F_i$ and suppose the family blocks. The group $\U_H$ acts
diagonally on $A^s$ by the power maps. The orbit of a tuple whose largest
coordinate conductor exponent is $t$ is a copy of $\U_t$. Indeed a
stabilizing element must be $1\bmod p^t$ from the coordinate of that exact
order, and this condition fixes the whole tuple. On the orbit the positions
where $F=0$ meet every translate of $J_t$ and thus number at least $\tau_t$.
The inverse image of any cover under reduction $\U_H\to\U_t$ is again a
cover, so $c_H\le c_t$.

Choose a positive unit representative $u$ for each class in $\U_H$ and put
$P(T)=\prod_{u\in\U_H}F([u]T)$. At each tuple as above at least
$(e_H/e_t)\tau_t\ge\tau_H$ factors vanish. Their product therefore has
multiplicity at least $\tau_H$ there. Apply Lemma~\ref{lem:grid} with
$g=\prod_{h\in S}g_h$; this polynomial has $G_r(g)=\sum_{h\in S}e_hr_h=N$.
Hence
\[
 e_H B\ \ge\ e_H G_r(F)=G_r(P)\ \ge\ N\tau_H .
\]
The same argument on $S\cap[1,t]$ proves each prefix bound.

Conversely, choose a minimal cover $D_h$ in each group, a primitive root
$\xi_h$ of order $p^h$, and let
\[
 F_*(T)=\prod_{h\in S}\prod_{u\in D_h}\bigl(T_1-(\xi_h^u-1)\bigr).
\]
The first coordinate of every candidate tuple has a power in the root set.
At zero this polynomial has valuation $\sum_h |D_h|r_h=\sum_h c_h$.
In one variable a blocking product must itself vanish at $\tau_h$ or more
distinct points of each shell. An integral series vanishing at $a$ of
positive valuation can be divided by $T-a$ with integral quotient, as
follows from summing $(T^n-a^n)/(T-a)$ termwise. Distinct roots persist after
each division. Every extracted factor has $G_r(T-a)=r_h$; this proves the
one-variable lower bound. All claimed equalities follow.
\end{proof}

\section{Sharp endpoints for the first two powers}\label{sec:two}
For $a\ge2$ prime to $p$ and $J=\{1,a,\ldots,a^{k-1}\}$, put
$d_h=\ord_{p^h}(a)$. Every orbit under multiplication by $a$ is a cycle of
length $d_h$. A cover must meet every cyclic run of $k$ consecutive
positions. Its minimum size on a cycle is $\ceil{d_h/k}$: each marked
position covers at most $k$ starts; taking positions $0,k,2k,\ldots<d_h$
works also across the seam. Thus
\[
 c_h=\frac{\ceil{d_h/k}}{d_h}.
\]
In particular, if $k\mid d_{h_0}$, every interval of $N$ shells beginning
at $h_0$ has exact obstruction cost $N/k$ in every dimension.

\begin{theorem}\label{th:end}
Let $p$ be odd, $J=\{1,2\}$, and $r\ge r_{h_0}$. Define $H_s(B)$
as the least endpoint $H\ge h_0$ such that every admissible family of
budget at most $B$ has a simultaneous nonzero specialization on
$(\bigcup_{h=h_0}^H Z_h)^s$. Let $d=\ord_p(2)$.
\begin{enumerate}
\item If $d$ is even, $H_s(B)=h_0+\lfloor 2B\rfloor$ for every $B\ge0$.
\item If $d$ is odd, $H_s(k/2)=h_0+k-1$ for every integer $k\ge1$,
and $H_s(0)=h_0$. For general $B$ with $\lfloor2B\rfloor\ge1$,
\[
 h_0+\lfloor 2B\rfloor-1\le H_s(B)\le h_0+\lfloor2B\rfloor.
\]
If $0\le B<1/2$ then $H_s(B)=h_0$.
\item In one dimension and for any $B$, the exact answer is the first $H$
such that $\sum_{h=h_0}^H\ceil{d_h/2}/d_h>B$.
\end{enumerate}
\end{theorem}
\begin{proof}
Write $\lambda=v_p(2^d-1)$. For odd $p$ the binomial theorem gives
$v_p((1+b)^p-1)=v_p(b)+1$ when $v_p(b)\ge1$, and exponentiation by an
integer prime to $p$ leaves this valuation unchanged. It follows that
$d_h=d\,p^{\max(0,h-\lambda)}$. Hence all orders have the parity of $d$.
The even case follows from exact cost $N/2$.

In the odd case $c_h=1/2+1/(2d_h)$, so $Nc_H>N/2$ on every nonempty window.
The possible extra cost in the explicit blocker is uniformly bounded since
\[
 \sum_{h\ge1}\frac1{d_h}=\frac{\lambda+1/(p-1)}{d}<1.
\]
Indeed $p^\lambda\le 2^d-1<p^d$ implies $\lambda\le d-1$.
Thus a window of $N$ shells can be blocked with strictly less than
$(N+1)/2$, and cannot be blocked at budget $N/2$. For $B=k/2$ the window
of $k$ shells always succeeds, whereas that of $k-1$ shells, if nonempty,
can be blocked. For arbitrary $B$, take $k=\lfloor2B\rfloor$. A window
of $k+1$ shells succeeds, while one of $k-1$ can be blocked. The remaining
assertions follow from Theorem~\ref{th:cover}.
\end{proof}

For the search range in the article, set $h_0$ to the first $h\ge1$ with
$r_h\le r/2$. Its sufficient endpoint for the first two powers is then
$h_0+\lfloor2B\rfloor$. At a half-integral positive budget, this endpoint
is optimal exactly when $\ord_p(2)$ is even. Both behaviors occur in
infinite families: primes $p\equiv7\pmod8$ have $2$ a square and
$(p-1)/2$ odd; at primes $p\equiv3\pmod8$, $2$ is a nonsquare and its
order is even. These statements follow from the supplementary law for $2$
in quadratic reciprocity and Dirichlet's theorem on progressions.

\begin{example}
At $p=7$, $d_h=3\cdot7^{h-1}$. The first shell costs $2/3$ in any
dimension. On the first two shells the exact one-variable cost is
$2/3+11/21=25/21$. In all dimensions a budget of $1$ already guarantees
success on these two shells: the orbit bound gives $2\cdot11/21>1$,
while the first shell can be blocked. At $p=5$ the first two shells
together cost exactly $1$ in all dimensions, and a third is necessary
at that budget.
\end{example}

\paragraph{Scope and remaining question.}
In several variables a tuple may have coordinates of different conductors;
its unit orbit is controlled by the largest of them. This is precisely
why the lower bound for mixed shells uses $c_t$ at a prefix endpoint.
For nonconstant densities it remains open here whether $\sum_{h\in S}c_h$
is the exact mixed-shell obstruction cost for every $s$. Powers divisible
by $p$ move between shells and require a different covering problem.
The argument uses no hypothesis on zeros of global $L$-functions.
Its contribution is the passage from a group covering on a shell to the
sharp analytic threshold in every dimension, and the resulting endpoint
classification for the first two powers.

\begin{thebibliography}{9}
\bibitem{proveit} V. Reshetnikov, \emph{ProveIt}, Generalized Riemann
Hypothesis article, sections \emph{Torsion specialization: selection theorems}
and \emph{Torsion specialization: exact costs, obstructions and coefficient
fields}, snapshot \texttt{4ccd0fdeda0b}.\enspace
\href{https://github.com/VladimirReshetnikov/ProveIt/blob/4ccd0fdeda0b1cc316d2dfdccfdbcfea9d4b77fb/Analysis/GeneralizedRiemannHypothesis/Article/sections/tors.tex}{Selection theorems}
\enspace and\enspace
\href{https://github.com/VladimirReshetnikov/ProveIt/blob/4ccd0fdeda0b1cc316d2dfdccfdbcfea9d4b77fb/Analysis/GeneralizedRiemannHypothesis/Article/sections/cost.tex}{exact costs}.
\bibitem{ballserra} S. Ball and O. Serra, \emph{Punctured combinatorial
Nullstellens\"atze}, Combinatorica \textbf{29} (2009), 511--522.
\href{https://doi.org/10.1007/s00493-009-2509-z}{doi:10.1007/s00493-009-2509-z}.
\bibitem{blackburn} S. R. Blackburn and J. F. McKee,
\emph{Constructing $k$-radius sequences}, Math. Comp. \textbf{81} (2012),
2439--2459. \href{https://arxiv.org/abs/1006.5812}{arXiv:1006.5812}.
\end{thebibliography}
\end{document}
